% Alternative wording of the paragraph below Table 5, restricted to the leak column
% (kept in case the two-column reading is ever reinstated). Not part of the build.

The remaining column of Table~\ref{tab:welfare-comparison} adds the labor side to the valuation. Every measure up to this point values the bundle alone; this one also accounts for the additional work that a programme requires -- so that a gain paid for in hours of labor no longer counts as a free gain. The wage is the natural price for that work, and it requires no preference parameter: at the calibrated baseline the household's marginal valuation of the last hour supplied equals the wage, so the leisure forgone and the wage income received offset each other to first order, and the correction is exact at that order with no elasticity to choose. Three patterns follow. First, wherever employment cannot move, the leisure term vanishes identically and the column reproduces $W_{CES^*}$ exactly; this is why the no-reallocation and the standard-CGE cells show no gap between the two, under any financing. Second, where employment does move, the leisure term takes the sign of that response: at the wage-pinned pole employment \emph{falls} under F1, so the leisure term registers a gain and the measure rises above the bundle value it qualifies ($+0.25$ to $+0.36$), while under external finance employment \emph{rises} by $1.78$~per cent and more than half of the $+3.56$ bundle gain turns out to be the value of the extra work, leaving $+1.59$. Third, the correction is expressed on the household's resources rather than on a consumption basket -- the one place in the table where the treatment of saving has to be stated.\footnote{Under the deferred convention saving counts as postponed consumption rather than as a leak, so the private block enters the bundle scaled by $1/(1-s)$ before valuation. That is not a re-scaling of the same object: it raises the private weight in the baseline bundle from $0.76$ to $0.79$ and the scale $V_0$ from $0.91$ to $1.00$, while leaving the per-sector relative changes untouched, so it changes what is valued rather than only how the correction is treated. The leak reading is reported throughout; the deferred variant is available on the same code path and is recorded in the replication materials.} We report the conservative reading, in which saving is a leak that never returns and only what the household consumes is valued. Treating it instead as postponed consumption is a robustness variant rather than a co-equal measure, since it changes the object valued -- re-weighting the baseline bundle toward private demand -- and not merely the treatment of the correction; it is not tabulated here, and the footnote records it.
